\documentclass[11pt, letterpaper]{article}
\usepackage[margin=1in]{geometry}
\usepackage{enumitem}
\usepackage{parskip}
\usepackage{titlesec}
\usepackage{booktabs}
\usepackage[hidelinks]{hyperref}

\titleformat{\section}{\normalfont\bfseries}{}{0em}{}[\titlerule]
\titlespacing{\section}{0pt}{8pt}{4pt}

\begin{document}

\begin{center}
    {\Large \textbf{Umbra: A Light and Shadow Puzzle Game}}\\[2pt]
    {\small CIS 5680 Group Game Project 1 --- Spring 2026}
\end{center}

\vspace{-4pt}

\section*{High Concept}

In \textbf{Umbra}, the world is split into two physical states---light and shadow---and only one
exists at a time in any given tile. Players manipulate light sources to recast shadows across the
environment, toggling regions between lit and dark to unlock new paths and solve spatial puzzles.

\section*{Genre}

Top-down 2D spatial puzzle game.

\section*{Main Features}

\begin{itemize}[leftmargin=1.5em, itemsep=3pt]

    \item \textbf{Light/Dark World Duality.} Every tile in the level belongs to either the light
    world or the shadow world at any given moment. Objects, platforms, and pathways exist
    \textit{only} in one world---a bridge visible in shadow vanishes when light floods in; a door
    only opens from the lit side. Rearranging light sources rewrites which half of the world is
    accessible.

    \item \textbf{Morning vs.\ Evening Toggle.} Players switch between a ``morning'' sun angle and
    an ``evening'' sun angle, flipping the direction and length of every shadow in the level
    simultaneously. This single binary toggle is the primary puzzle lever: a path blocked in
    morning becomes open in evening, and vice versa. Mastering \textit{when} to toggle is the
    core skill.

    \item \textbf{Shadows as Physical Geometry.} Shadow regions are not merely visual---they are
    solid, traversable terrain. A chasm that cannot be crossed under full light becomes passable
    once an object is repositioned to cast a shadow-bridge across it. Moving that object, or
    toggling the sun angle, destroys the bridge. Players must plan object placement and sun state
    together.

    \item \textbf{Interactive Light Sources.} Portable lanterns, pivotable wall-mounted spotlights,
    and moveable blocking objects (pillars, crates) let players sculpt the shadow layout before
    committing to a move. Each level is a static logic puzzle: find the one arrangement that opens
    a continuous path from start to exit across both light and shadow tiles.

    \item \textbf{Minimalist Silhouette Aesthetic.} Clean geometric shapes with stark contrast
    between amber-lit tiles and cool blue-grey shadow tiles make the current world state instantly
    legible. The visual language is simple enough that the puzzle logic---not the art---demands
    the player's attention.

\end{itemize}

\section*{Player Challenge / Motivation}

The player is an explorer navigating a world where light and shadow are as solid as stone. Each
self-contained level is a spatial logic puzzle: determine the correct positions of light sources and
the correct sun state (morning or evening) to open a continuous walkable path to the exit. There is
no time pressure and no combat---only the satisfaction of reading the environment and finding the
solution hidden inside it.

\section*{Design Goals}

\begin{itemize}[leftmargin=1.5em, itemsep=2pt]
    \item \textbf{Elegant.} Two rules---shadows are solid, and the sun has two states---combine to
    produce a surprisingly large puzzle space without requiring the player to learn many systems.
    \item \textbf{Satisfying.} Every level has a clean ``aha'' moment when the correct light
    arrangement snaps into place and the path becomes obvious. Solutions should feel discovered,
    not stumbled upon.
    \item \textbf{Readable.} The player should always be able to see clearly which tiles are lit
    and which are in shadow, and understand exactly what will change when they move a light source
    or toggle the sun state. No hidden information.
\end{itemize}

\section*{Target Customer}

Fans of pure puzzle games such as \textit{The Witness}, \textit{Portal}, and \textit{Stephen's
Sausage Roll}---players who want a single elegant mechanic taken to its logical depth rather than
narrative or action content. Accessible to ages 12 and up; pick-up-and-play sessions of 5--15
minutes per level.

\section*{Unique Selling Points}

\begin{itemize}[leftmargin=1.5em, itemsep=2pt]
    \item The world is literally divided into two physical halves by light and shadow; the player
    inhabits both by controlling where the boundary falls.
    \item A single morning/evening sun-angle toggle simultaneously remaps every shadow in the
    level, creating a whole new traversable topology from the same geometry.
    \item Pure puzzle focus with no combat, no timers, and no randomness---every solution is
    deterministic and elegant.
\end{itemize}

\section*{Competition}

\textit{Contrast} (2013) projects a character's shadow onto walls as a 2D platforming layer, but
is story-driven and does not feature a light/dark world split. \textit{The Witness} and
\textit{Portal} share the clean, mechanic-first puzzle philosophy but use completely different
core mechanics. \textit{Umbra} is unique in using directional shadow casting as the sole map
editor and spatial puzzle variable in a top-down environment.

\section*{Target Hardware}

PC (Windows), implemented in Unreal Engine 5. Mouse-and-keyboard or gamepad supported.

\end{document}
